\scanpage{186}
% archive sheet number: 185

\deleted{\unsure{cogèbres}} \deleted{groupes} objets
\[
\text{(ind (algébriques finis))}\simeq\text{coalgèbres comm \unsure{nn. counitaires}}
\]
\margin{car \deleted{\unsure{on}} on est sur un corps, et on peut utiliser la dualité}
\deleted{\unsure{les}} \uline{engendrées par} des \uline{coalgèbres de} \uline{rang fini sur $k$}

\deleted{[une telle coalgèbre est en effet engendrée par des coalgèbres}
(cette dernière condition sans doute superflue, i.e. une multiplication \inserted{\unsure{comm non unitaire}} \uline{continue} sur un $k^I$ \deleted{\unsure{n'est}} \inserted{\unsure{n'impose}} pas nécessairement une \deleted{\unsure{système}} \unsure{limite} projective d'algèbres de rang fini sur $k$)

\margin{$\Delta(u)\in U\otimes U)$}
\diagram{\begin{tikzcd}[ampersand replacement=\&]
U \arrow[dr,leftrightarrow] \& \unsure{V|U} \arrow[dl,leftrightarrow] \\
V^* \& U^*
\end{tikzcd}}

\deleted{\unsure{gr}}
\[
\underbrace{\text{groupe (ind (algébriques finis))}}_{\text{groupe formel}}
\simeq\text{coalgèbre du type précédent,}
\]
munie d'une loi d'algèbre ass. commutative unitaire compatible avec la précédente

\margin{\uline{Exemple} l'Algèbre affine d'un groupe algébrique linéaire ou celle d'un pro. (groupe algébrique linéaire)}

d'après Cartier, \inserted{en car. $p$} une telle algèbre serait limite inductive d'algèbres de \inserted{\unsure{ce type, qui sont}} de type fini du pt de vue associatif [du moins sur un corps alg. clos, où la partie discrète se sépare de la partie infinitésimale]

Donc il en résulterait une dualité entre
\[
\mathrm{pro}\ (\text{grps algébriques affines})\sim\mathrm{grp}\ (\text{ind. Alg. finis}).
\]
Le plus beau serait maintenant qu'on ait le résultat que cel\unsure{a} correspond aux isomorphismes \uline{\unsure{algébriques}}
